\section{Conclusion}
\label{sec:conclusion}
The present implementation demonstrates a working proof-of-concept generative-design pipeline, a reduced freeform-object experiment, and a public VSP Airshow smoke-training run that can be executed in the installed environment. The Airshow addition moves the paper beyond synthetic-only training evidence by converting 355 traceable public geometry records into a manifest-backed \(16^3\) voxel corpus. It also exposes the current generator's limit: the three Airshow-checkpoint flight-path samples are nonempty and solver-runnable, but all fail the aircraft-specific \texttt{span\_sanity} screen.

The OpenFOAM cross-check completes on the shared centered-cube validation object and gives the repository an external CFD comparison path. The reduced protocol also assembles a passing package-level evidence report with manifest validation, heuristic aircraft-shape checks, grounded condition-response sweeps, manufacturing-bounds checks, and finite baseline statistics. The next revision should therefore prioritize generated-sample validity, solver calibration, and larger ablation studies before presenting stronger aerodynamic conclusions.

Most importantly, the repository should not yet be described as a fully AI-driven airplane generator conditioned on flight profile, manufacturing method, or structural requirements. What exists today is structured-conditioning plumbing plus a reduced evidence package and a grounded public-corpus smoke run. That stronger version needs more evidence first: higher-resolution or validity-bearing aircraft-like training runs, conditional ablations, stricter structural and aerodynamic checks, and repeated runs beyond the present reduced local protocol.
